% ============================================================
%  Simba Mini Banking System – Technical Report
%  TEMB 1202: Computing II | Kyambogo University | Group Two
% ============================================================
\documentclass[12pt, a4paper]{article}

% ---------- Packages ----------
\usepackage[T1]{fontenc}
\usepackage[utf8]{inputenc}
\usepackage{lmodern}
\usepackage[margin=2.5cm]{geometry}
\usepackage{graphicx}
\usepackage{xcolor}
\usepackage{booktabs}
\usepackage{array}
\usepackage{longtable}
\usepackage{listings}
\usepackage{fancyhdr}
\usepackage{titlesec}
\usepackage{hyperref}
\usepackage{setspace}
\usepackage{parskip}
\usepackage{microtype}
\usepackage{enumitem}
\usepackage{mdframed}

% ---------- Colour palette (black & white only) ----------
\definecolor{codeGray}{RGB}{245, 245, 245}
\definecolor{commentGreen}{RGB}{80, 80, 80}
\definecolor{keywordBlue}{RGB}{0, 0, 0}
\definecolor{stringRed}{RGB}{0, 0, 0}

% ---------- Listings (C code) ----------
\lstdefinestyle{cstyle}{
    backgroundcolor=\color{codeGray},
    basicstyle=\ttfamily\small,
    keywordstyle=\color{black}\bfseries,
    commentstyle=\color{black}\itshape,
    stringstyle=\color{black},
    numberstyle=\tiny\color{black},
    breaklines=true,
    breakatwhitespace=false,
    numbers=left,
    numbersep=8pt,
    frame=single,
    rulecolor=\color{black},
    tabsize=4,
    showstringspaces=false,
    captionpos=b,
    language=C
}
\lstset{style=cstyle}

% ---------- Section formatting ----------
\titleformat{\section}
    {\Large\bfseries\color{black}}
    {\thesection.}{0.6em}{}[\titlerule]

\titleformat{\subsection}
    {\large\bfseries\color{black}}
    {\thesubsection}{0.5em}{}

\titleformat{\subsubsection}
    {\normalsize\bfseries\color{black}}
    {\thesubsubsection}{0.4em}{}

% ---------- Header / Footer ----------
\pagestyle{fancy}
\fancyhf{}
\fancyhead[L]{\small\color{black}\textbf{TEMB 1202 – Computing II}}
\fancyhead[R]{\small\color{black}Simba Mini Banking System}
\fancyfoot[C]{\small\color{black}Page \thepage}
\renewcommand{\headrulewidth}{0.4pt}
\renewcommand{\footrulewidth}{0.4pt}

% ---------- Hyperref ----------
\hypersetup{
    colorlinks=true,
    linkcolor=black,
    urlcolor=black,
    pdftitle={Simba Mini Banking System – Technical Report},
    pdfauthor={Group Two, Kyambogo University}
}

% ---------- Spacing ----------
\setstretch{1.15}
\setlength{\parindent}{0pt}
\setlength{\parskip}{6pt}

% ============================================================
\begin{document}

% ======================== TITLE PAGE ========================
\begin{titlepage}
    \centering
    \vspace*{1.5cm}

    {\Large\bfseries\color{black} Kyambogo University}\\[0.3cm]
    {\normalsize Department of Biomedical and Mechatronics Engineering}\\[0.1cm]
    {\normalsize TEMB 1202: Computing II Assignment}

    \vspace{1.2cm}
    \textcolor{black}{\rule{12cm}{2pt}}
    \vspace{0.5cm}

    {\Huge\bfseries\color{black} SIMBA MINI}\\[0.3cm]
    {\Huge\bfseries\color{black} BANKING SYSTEM}

    \vspace{0.4cm}
    {\large\itshape Technical Report \& Documentation}

    \textcolor{black}{\rule{12cm}{2pt}}

    \vspace{1.5cm}

    \begin{tabular}{rl}
        \textbf{Group}      & Two \\[4pt]
        \textbf{Course}     & TEMB 1202 – Computing II \\[4pt]
        \textbf{Assignment} & Mini Banking System (Group 2) \\[4pt]
        \textbf{Language}   & C (Standard C99) \\[4pt]
        \textbf{File}       & \texttt{mini\_banking\_system.c} \\[4pt]
        \textbf{Source Code}& \href{https://github.com/MackonAjax/MiniBankingSystem/blob/main/main.c}{\texttt{github.com/MackonAjax/MiniBankingSystem}} \\
    \end{tabular}

    \vfill
    {\small\color{black} Kyambogo University $\cdot$ Department of Biomedical and Mechatronics Engineering}
\end{titlepage}

% ======================== TOC ========================
\tableofcontents
\thispagestyle{fancy}
\newpage

% ======================== SECTIONS ========================

\section{Introduction / Problem Statement}

The \textbf{Simba Mini Banking System} is a menu-driven C program developed as part of
the TEMB 1202 Computing~II course assignment at Kyambogo University. The assignment
required Group Two to design and implement a fully functional banking simulation that
demonstrates core C programming concepts including functions, file I/O, loops, conditional
statements, arrays, and string handling. The full source code is available on GitHub at
\href{https://github.com/MackonAjax/MiniBankingSystem/blob/main/main.c}{\texttt{github.com/MackonAjax/MiniBankingSystem}}.

The central problem addressed is the lack of a lightweight, terminal-based tool that
simulates real-world banking operations — account creation, authentication, money
deposits, withdrawals, balance tracking, and account management — entirely in C without
reliance on external databases or graphical interfaces.

The program stores user account data persistently in a flat-file database
(\texttt{simba\_users.db}), allowing information to survive between program sessions. It
enforces business rules such as minimum opening balance, withdrawal limits, and
low-balance warnings, giving it the feel of a real banking application.

% -------------------------------------------------------
\section{System Features}

\begin{description}[leftmargin=2em, labelwidth=!, labelindent=0pt,
                    font=\bfseries\color{black}]

    \item[Account Registration (Sign Up)]
        New users can create a Simba Bank account by providing their full name, phone
        number, and a password. The system validates all inputs and requires a minimum
        opening deposit of UGX~10{,}000 before the account is created.

    \item[Authentication (Log In)]
        Existing users authenticate using their registered phone number and password.
        The system reads from the persistent flat-file database to verify credentials,
        rejecting unrecognised or incorrect details.

    \item[Deposit Money]
        Authenticated users may deposit any positive amount into their account. A
        transaction confirmation message is displayed after every successful deposit.

    \item[Withdraw Money]
        Withdrawals are processed with checks: the balance must not drop below the
        minimum required threshold, and a warning is displayed for unusually large
        withdrawal amounts.

    \item[Send Money]
        Users can transfer funds to another registered Simba Bank account using the
        recipient's phone number, subject to the same balance constraints as
        withdrawals.

    \item[View Transaction History]
        The dashboard provides an option to display a summary of past transactions
        performed on the account during the current session.

    \item[Account Information Update]
        Users can update their stored account details (name, phone number, password)
        from within the authenticated dashboard.

    \item[Account Deletion]
        Users may permanently close and delete their account from the system.

    \item[Balance Advisory]
        Whenever the account balance falls at or below UGX~15{,}000, the system
        automatically displays an advisory message encouraging the user to make a
        deposit.

    \item[Persistent File Storage]
        All account records are saved to \texttt{simba\_users.db} using C file~I/O,
        ensuring data is retained between program runs.

    \item[Continuous Operation]
        The application runs in a continuous loop until the user explicitly selects the
        \textit{Quit} option from the main entry menu.

\end{description}

% -------------------------------------------------------
\section{Program Design and Structure}

\subsection{Architecture Overview}

The program follows a \textbf{modular, function-based architecture}. The \texttt{main()}
function acts solely as an application controller, running a loop that calls
\texttt{entry\_menu()} and exits when a sentinel value (\texttt{10}) is returned. All
business logic is encapsulated in dedicated functions, improving readability and
maintainability.

\subsection{Module / Function Map}

\begin{longtable}{>{\bfseries\color{black}}p{3.5cm}
                  >{\ttfamily}p{5cm}
                  p{6cm}}
    \toprule
    \normalfont\bfseries Module & \normalfont\bfseries Functions & \normalfont\bfseries Responsibility \\
    \midrule
    \endfirsthead
    \toprule
    \normalfont\bfseries Module & \normalfont\bfseries Functions & \normalfont\bfseries Responsibility \\
    \midrule
    \endhead
    \bottomrule
    \endfoot

    Entry \& Navigation   & entry\_menu()
                          & Main menu loop; routes user to sign-up, login, help, or quit \\[6pt]

    Authentication        & sign\_up()\newline log\_in()
                          & Account creation and credential verification \\[6pt]

    Validation            & validate\_sign\_up\_details()
                          & Input sanity-checks for name, phone, and password \\[6pt]

    File Persistence      & write\_user\_to\_file()
                          & Appends new account records to \texttt{simba\_users.db} \\[6pt]

    User Dashboard        & users\_dashboard()
                          & Post-login hub: deposit, withdraw, send, history, logout, delete \\[6pt]

    Transactions          & deposit()\newline withdraw()\newline send\_money()
                          & Core financial operations with rule enforcement \\[6pt]

    Account Management    & delete\_account()\newline change\_account\_information()
                          & Account deletion and info updates \\

\end{longtable}

\subsection{Data Flow}

When the program starts, \texttt{main()} calls \texttt{entry\_menu()} in a loop. Selecting
option~1 triggers \texttt{sign\_up()}, which collects input, calls
\texttt{validate\_sign\_up\_details()}, and — on success — calls
\texttt{write\_user\_to\_file()} to persist the record, then immediately forwards the session
to \texttt{users\_dashboard()}. Selecting option~2 invokes \texttt{log\_in()}, which reads
\texttt{simba\_users.db} to authenticate the user before loading the dashboard. All
subsequent operations (deposit, withdraw, send, etc.) are handled from within the
dashboard loop.

\subsection{File Storage Format}

Account records are stored one per line in \texttt{simba\_users.db} using a
\textbf{pipe-delimited format}:

\begin{mdframed}[backgroundcolor=codeGray, linecolor=black,
                 linewidth=0.8pt, innerleftmargin=10pt, innerrightmargin=10pt,
                 innertopmargin=8pt, innerbottommargin=8pt]
\texttt{phone\_number | name | password | account\_balance}\\
\texttt{e.g.\ \ 0712345678 | John Doe | pass123 | 50000}
\end{mdframed}

The file is opened in append mode (\texttt{"a"}) during registration, ensuring existing
records are never overwritten.

% -------------------------------------------------------
\section{Important Functions Used}

\subsection{\texttt{entry\_menu()}}

\textbf{Return type:} \texttt{int} \quad \textbf{Parameters:} none

The gateway function of the entire application. It renders the main menu and dispatches
the user's selection to the appropriate handler. Returning the value \texttt{10} signals the
main loop to terminate. Invalid integer input triggers a recursive call to redisplay the menu.

\begin{lstlisting}[caption={Skeleton of \texttt{entry\_menu()}}]
int entry_menu(){
    // Displays main menu: Create Account | Login | Help | Quit
    // Returns 10 to break the main application loop on Quit
    ...
    if(option == 4){ return 10; }
    ...
}
\end{lstlisting}

\subsection{\texttt{sign\_up()}}

\textbf{Return type:} \texttt{void} \quad \textbf{Parameters:} none

Manages the full account-creation flow. It flushes the input buffer after \texttt{scanf()} to
prevent stale newline characters from corrupting \texttt{fgets()} calls, collects
name/phone/password via \texttt{fgets()}, delegates validation, prompts for an initial
deposit (minimum UGX~10{,}000 enforced with a \texttt{goto} retry loop), persists the
record, and launches the user dashboard.

\begin{lstlisting}[caption={Pseudocode outline of \texttt{sign\_up()}}]
void sign_up(){
    // 1. Flush \n from stdin
    // 2. Collect name, phone_number, password via fgets()
    // 3. Call validate_sign_up_details() -- retry on failure
    // 4. Enforce minimum deposit >= 10,000 (goto retry)
    // 5. write_user_to_file() then users_dashboard()
}
\end{lstlisting}

\subsection{\texttt{validate\_sign\_up\_details()}}

\textbf{Return type:} \texttt{int} \quad
\textbf{Parameters:} \texttt{name[30]}, \texttt{phone\_number[11]}, \texttt{password[30]}

A dedicated validation layer that checks each input field independently and returns a
numeric error code:
\begin{itemize}
    \item \textbf{0} — all valid
    \item \textbf{1} — name error
    \item \textbf{2} — phone number error
    \item \textbf{3} — password error (must be 6–30 characters)
\end{itemize}
This separation of concerns keeps \texttt{sign\_up()} clean and makes future validation
rule changes straightforward.

\subsection{\texttt{write\_user\_to\_file()}}

\textbf{Return type:} \texttt{void} \quad
\textbf{Parameters:} \texttt{name}, \texttt{phone\_number}, \texttt{password},
\texttt{account\_balance}

Responsible for persisting a new account to disk. Before writing, it strips trailing newline
characters from all string fields using \texttt{strcspn()} — a critical step because
\texttt{fgets()} includes the newline in the captured string. The record is appended with
\texttt{fprintf()} in the pipe-delimited format described in Section~\ref{sec:design}.

\begin{lstlisting}[caption={Implementation of \texttt{write\_user\_to\_file()}}]
void write_user_to_file(char name[30], char phone_number[11],
                        char password[30], int account_balance){
    FILE *file_ptr = fopen("./simba_users.db", "a");
    name[strcspn(name, "\n")]               = '\0'; // strip newline
    phone_number[strcspn(phone_number,"\n")]= '\0';
    password[strcspn(password, "\n")]       = '\0';
    fprintf(file_ptr, "%s | %s | %s | %d\n",
            phone_number, name, password, account_balance);
    fclose(file_ptr);
}
\end{lstlisting}

\subsection{\texttt{users\_dashboard()}}

\textbf{Return type:} \texttt{void} \quad
\textbf{Parameters:} \texttt{name}, \texttt{phone\_number}, \texttt{password},
\texttt{account\_balance}

The post-authentication control hub. It loops continuously, displaying the account holder's
name, phone number, current balance, and an advisory message if the balance is at or
below UGX~15{,}000. The user chooses from six options:

\begin{enumerate}
    \item Deposit
    \item Send Money
    \item Withdraw
    \item View History
    \item Log Out
    \item Delete Account
\end{enumerate}

Selecting \textit{Log Out} breaks the inner loop, returning control to \texttt{entry\_menu()}.

\subsection{\texttt{log\_in()}}

\textbf{Return type:} \texttt{void} \quad \textbf{Parameters:} none

Collects the user's phone number and password, then searches \texttt{simba\_users.db}
for a matching record. On success it loads \texttt{users\_dashboard()} with the retrieved
account data.

\begin{mdframed}[backgroundcolor=yellow!12, linecolor=black,
                 linewidth=1pt, innerleftmargin=10pt, innerrightmargin=10pt,
                 innertopmargin=8pt, innerbottommargin=8pt]
\textbf{Note:} A syntax error (missing semicolon after a \texttt{printf} call) and a string
comparison bug (using \texttt{==} instead of \texttt{strcmp()}) are present in the current
draft and are flagged under Section~6 (Challenges).
\end{mdframed}

% -------------------------------------------------------
\section{Design Decisions and Justifications}
\label{sec:design}

\subsection{Flat-file Database (\texttt{simba\_users.db})}

A plain-text, pipe-delimited file was chosen for persistence because it requires no external
libraries, is human-readable for debugging, and is entirely within the scope of C standard
I/O. The append mode (\texttt{'a'}) guarantees that adding a new account never corrupts
existing records.

\subsection{\texttt{fgets()} over \texttt{scanf()} for String Input}

\texttt{scanf()} with \texttt{\%s} stops at whitespace, making it unsuitable for capturing
full names. \texttt{fgets()} reads entire lines including spaces, providing a more robust
input method. A deliberate buffer-flush loop
(\texttt{while~(getchar()~!=~'\textbackslash n')}) is used before \texttt{fgets()} calls to
discard the newline left by prior \texttt{scanf()} calls.

\subsection{Modular Function Architecture}

Every logical action — validation, file writing, menu display, dashboard — is encapsulated
in its own function. This matches the assignment's requirement for user-defined functions
and makes individual components independently testable and replaceable.

\subsection{Numeric Sentinel Return (\texttt{10}) from \texttt{entry\_menu()}}

Instead of a global flag variable, a specific return value (\texttt{10}) communicates the
exit intent up to \texttt{main()}. This keeps side effects minimal and makes the control
flow explicit and easy to trace.

\subsection{\texttt{goto} for Retry Logic in \texttt{sign\_up()}}

A labelled \texttt{goto} is used to retry the initial deposit prompt when the user enters a
value below the minimum. While \texttt{goto} is generally discouraged in modern C, its use
here is localised, clearly labelled (\texttt{account\_creation:}), and avoids introducing an
additional nested function or loop just for a single validation step.

\subsection{Minimum Opening Balance of UGX~10{,}000}

A minimum opening deposit enforces the assignment rule that \textit{``a new account must
start with a minimum opening balance,''} preventing zero-balance ghost accounts and
mirroring real banking practice.

\subsection{Low-Balance Advisory at UGX~15{,}000}

The threshold is set deliberately above the 10{,}000 minimum so that users receive a
warning before reaching the absolute floor, giving them time to act. This reflects the
assignment rule: \textit{``if the balance becomes too low, advise the user to deposit
money.''}

\subsection{\texttt{strcspn()} for Newline Stripping}

\texttt{fgets()} captures the trailing newline as part of the string. Using
\texttt{strcspn(str,~"\textbackslash n")} to locate and null-terminate at that position is a
clean, standard-library approach that avoids manual iteration and is less error-prone than
alternatives like \texttt{strtok()}.

% -------------------------------------------------------
\section{Challenges Encountered}

\subsection{Input Buffer Contamination}

Mixed use of \texttt{scanf()} (for integers) and \texttt{fgets()} (for strings) caused the
newline character left in \texttt{stdin} by \texttt{scanf()} to be immediately consumed by
the subsequent \texttt{fgets()} call, resulting in empty string captures. This was resolved
by inserting a buffer-flush loop:

\begin{lstlisting}[caption={Buffer-flush idiom}]
while ((c = getchar()) != '\n' && c != EOF);
\end{lstlisting}

between the two input methods wherever they appear together.

\subsection{Trailing Newlines in File Records}

\texttt{fgets()} includes the newline at the end of captured strings. Without stripping
these, the database records would contain embedded newlines, breaking the pipe-delimited
format and causing comparison failures during login. \texttt{strcspn()} was applied to all
three string fields before writing to file.

\subsection{String Comparison in \texttt{log\_in()}}

The draft \texttt{log\_in()} function contains the line:

\begin{lstlisting}[caption={Bug: pointer comparison instead of string comparison}]
if(phone_number == "q"){ return; }  // BUG: compares pointer address
\end{lstlisting}

This compares a pointer address rather than string contents — always evaluating false in
C. The correct approach is:

\begin{lstlisting}[caption={Fix: use \texttt{strcmp()}}]
if(strcmp(phone_number, "q") == 0){ return; }  // CORRECT
\end{lstlisting}

This is a known bug in the current version and is scheduled for correction.

\subsection{Missing Semicolon in \texttt{log\_in()}}

A syntax error — a missing semicolon after a \texttt{printf()} call in \texttt{log\_in()} —
prevents the current code from compiling as written. This was identified during code review
and will be patched before the final submission.

\subsection{Git Merge Conflict in Source File}

The repository history shows an unresolved Git merge conflict in the source file
(\texttt{<<<<<<< Updated upstream} / \texttt{>>>>>>> Stashed changes} markers). The
conflict arose when team members developed separate branches simultaneously without
coordinating a merge. The conflict was manually resolved by retaining the more complete
``Stashed changes'' version of \texttt{main()}.

\subsection{Recursive \texttt{sign\_up()} Calls on Validation Failure}

Early iterations used recursive calls to \texttt{sign\_up()} on validation failure, which
risks stack overflow if a user enters invalid data many times in succession. A future
improvement would replace recursion with an iterative loop within \texttt{sign\_up()}
itself.

\subsection{Incomplete Transaction Functions}

The \texttt{deposit()}, \texttt{withdraw()}, and \texttt{send\_money()} function bodies
remain as stubs in the current version. Their logic — updating the in-memory balance and
rewriting the file record — is planned but not yet implemented. This is the primary remaining
development task before the submission deadline.

% -------------------------------------------------------
\section{Conclusion}

The \textbf{Simba Mini Banking System} successfully demonstrates the application of core C
programming concepts in a practical, real-world context. The program achieves its primary
objectives: it provides a continuous menu-driven interface, enforces simple authentication,
supports account creation with validation and minimum-balance rules, and persists data
across sessions using file I/O. The complete source code is hosted publicly at
\href{https://github.com/MackonAjax/MiniBankingSystem/blob/main/main.c}{\texttt{github.com/MackonAjax/MiniBankingSystem}}.

The modular design — with clearly named, single-responsibility functions — makes the
codebase readable, maintainable, and extensible. The deliberate separation of validation
logic (\texttt{validate\_sign\_up\_details()}), persistence
(\texttt{write\_user\_to\_file()}), and presentation (\texttt{users\_dashboard()}) reflects
sound software engineering principles even within the constraints of an introductory C
assignment.

A number of known issues — the string comparison bug in \texttt{log\_in()}, the missing
semicolon, and the incomplete transaction function bodies — have been identified and
documented. These will be resolved before the final submission. The team also intends to
replace recursive retry logic with iterative loops to improve robustness.

Overall, the project provided valuable hands-on experience with practical C challenges
such as buffer management, file~I/O, input validation, and collaborative version control.
The lessons learned — particularly around mixed input methods and Git merge discipline —
will directly inform better practices in future development work.

\vspace{1.5cm}
\begin{center}
    \textcolor{black}{\rule{10cm}{1pt}}\\[0.4cm]
    {\small\itshape Group Two $\;\cdot\;$ TEMB 1202: Computing II $\;\cdot\;$
    Kyambogo University}
\end{center}

\end{document}
